\section{The dissipative vortex-layer programme: exact scales and residuals}
\label{sec:forced-program-scales}

The four-page statement points to the programme of Diego C\'ordoba and Luis
Mart\'inez-Zoroa.  A primary dissipative predecessor is their joint paper with
Fan Zheng, \emph{Finite time blow-up for the hypodissipative Navier--Stokes
equations with a force in $L^1_tC_x^{1,\epsilon}\cap L^\infty_tL_x^2$},
\href{https://arxiv.org/abs/2407.06776v2}{arXiv:2407.06776v2}, and its published
version, \href{https://doi.org/10.1007/s00205-026-02198-0}{Archive for Rational
Mechanics and Analysis 250 (2026), article 38}.
Their operator is
\[
 \widehat{\Lambda^\alpha v}(\xi)=|\xi|^\alpha\widehat v(\xi),\qquad
 \Lambda^\alpha=(-\Delta)^{\alpha/2}.
\]
In this convention the classical viscosity is $\nu\Lambda^2=-\nu\Delta$.
The theorem stated in those sources concerns
$0\leq\alpha<(22-8\sqrt7)/9$ and a force in the indicated time-integrated
H\"older class. It does not state smoothness of the force through the terminal
time. The following calculations reproduce the actual scale mechanism and
repair specified transcription and residual identities. They do not reproduce
the source's analytic flow-map and singular-integral estimates in full.
Whenever an estimate from that paper is used below, it is identified as a
source estimate; the algebra applied to its explicit right-hand side is proved
here. In particular this section asserts no new blow-up theorem.

\subsection{All scale variables and the source's initial four constraints}

Retain $N>1$, $R>1$, and integer $n\geq1$. Define
\[
 M_n=N^{R^n},\quad A=N^a,\quad A_n=A^{R^n}=M_n^a,
 \quad L=N^\ell,\quad L_n=L^{R^n}=M_n^\ell.
\]
The index $n$ is never substituted for $R^n$ in these definitions. The
parameters $a,\ell$ are exponents, while $A_n,L_n,M_n$ are positive magnitudes.
Let $K_n$ denote the stretching factor of layer $n$.
The source's introductory scale calculation associates the four quantities
\[
 \frac{M_n^\alpha}{A_{n-1}},\qquad
 \frac{A_nM_n^s}{K_n},\qquad
 \frac{K_nM_{n-1}}{L_n},\qquad
 \frac{A_n^2L_nM_n^{s-1}}{A_{n-1}}
 \tag{FS1}\label{eq:fs-four}
\]
with dissipation versus stretching, force needed to start a layer, spatial
support within the outer wavelength, and integrated quadratic interaction,
respectively. Their being at most one is a model of those four costs, not an
identity for the full velocity field.

Here is the exact algebra of that model. The second and third inequalities give
$K_n\geq A_nM_n^s$ and $L_n\geq A_nM_n^sM_{n-1}$. Substituting this lower
bound into the fourth gives
\[
 A_n^3M_n^{2s-1}M_{n-1}\leq A_{n-1}.
\]
On setting $A_{n-1}=M_n^\alpha$ and therefore
$A_n=M_{n+1}^\alpha$, the left side divided by the right side is
\[
 M_{n+1}^{3\alpha}M_n^{2s-1-\alpha}M_{n-1}
 =N^{R^{n-1}[3\alpha R^2+(2s-1-\alpha)R+1]}.
\]
Since $N>1$, its being at most one is exactly
\[
 3\alpha R^2+(2s-1-\alpha)R+1\leq0.                 \tag{FS2}
\]
For $\alpha>0$, this quadratic can be nonpositive at positive $R$ only if
$2s-1-\alpha<0$ and $(2s-1-\alpha)^2\geq12\alpha$.
The latter two requirements together are exactly
$s\leq(1+\alpha)/2-\sqrt{3\alpha}$. Indeed, the product of the two roots is
$1/(3\alpha)>0$; their positive sum is equivalent to the stated sign condition.
This is the introductory calculation. The full source chooses a smaller
regularity margin to control additional interactions, as follows.

\subsection{The precise additional exponent that produces the cutoff}

For the rest of this section set
\begin{align}
 \alpha_0&=\frac{22-8\sqrt7}{9},&0&<\alpha<\alpha_0,\nonumber\\
 R&=\sqrt{\frac{2}{7\alpha}},&
 a&=\sqrt{\frac{2\alpha}{7}},\nonumber\\
 s&=\frac{2+3\alpha-2\sqrt{14\alpha}}4,&
 \ell&=1+\alpha-s-2a.                              \label{eq:fs-parameters}
\end{align}
These are the parameters in Section 4 of the source. In particular
\[
 a=\alpha R,\quad \frac1R=\frac72a,\quad
 \sqrt{14\alpha}=7a,\quad
 A_{n-1}=M_n^\alpha,\quad
 \ell=\frac{2+\alpha}{4}+\frac32a.                \tag{FS3}\label{eq:fs-id}
\]
Each identity follows by squaring the positive square roots, or substituting
$4s=2+3\alpha-14a$; no coordinate or physical coefficient has been changed.

The source's outer-velocity error, after integrating its duration, has the
explicit scale
\[
 A_nM_n^{r+1}L_n^{-(3-\delta)}N^{2R^{n-1}}\log M_n.
                                                        \tag{FS4}\label{eq:fs-outer}
\]
For arbitrary $a=\alpha R$ and $\ell=1+\alpha-s-2a$, its exponent at
$\delta=0$ is exactly
\[
 a+r+1-3\ell+\frac2R
 =7\alpha R+\frac2R-2-3\alpha+r+3s.
\]
At $r=s$, the $R$-dependent part satisfies the identity
\[
 7\alpha R+\frac2R-2\sqrt{14\alpha}
 =\left(\sqrt{7\alpha R}-\sqrt{\frac2R}\right)^2\geq0.
\]
Equality holds exactly at the value of $R$ in \eqref{eq:fs-parameters}.
Thus $s$ in that equation is exactly the value at which the optimized
outer-error exponent at $r=s$ is zero. For the actual $r<s$ it becomes $r-s$.
This explains the numerical cutoff directly from the extra interaction cost.

We prove the admissible range without decimal approximations. The inequality
$s>0$ is $2+3\alpha>2\sqrt{14\alpha}$. Both sides are positive, so squaring is
equivalent and gives
\[
 9\alpha^2-44\alpha+4>0,
 \qquad
 9\alpha^2-44\alpha+4
 =9(\alpha-\alpha_0)
   \left(\alpha-\frac{22+8\sqrt7}{9}\right).
\]
The scale requirement $R>1$ implies $\alpha<2/7$. In this interval the second
root exceeds $2/7$, so $s>0$ is precisely $\alpha<\alpha_0$.
Furthermore $8\sqrt7>21$ because $448>441$, giving $\alpha_0<1/9$.
The following inequalities and identities will be used without losing their
strict margins:
\begin{gather}
 0<s<\frac12,\quad 0<a<\frac15,\quad a>\alpha,
 \quad 0<\ell<\frac56,\quad s+a<1,                 \label{eq:fs-margins}\\
 \ell-a-\frac1R-s=\frac{a-\alpha}{2}>0,\qquad
 3\ell-1-\frac2R=s+a>0.                           \label{eq:fs-separation}
\end{gather}
For $s<1/2$, square $3\alpha<2\sqrt{14\alpha}$, obtaining
$9\alpha<56$, valid for $\alpha<1/9$. Next
$a^2=2\alpha/7<2/63<1/25$, and $a>\alpha$ follows from
$\alpha<2/7$. The displayed alternative expression for $\ell$ makes it
positive and gives
\[
 \ell<\frac{19}{36}+\frac3{10}
 =\frac{149}{180}<\frac56.
\]
Also $s+a=(2+3\alpha-10a)/4<7/12<1$.
Finally direct substitution of $4s=2+3\alpha-14a$ and $1/R=7a/2$ proves both
identities in \eqref{eq:fs-separation}. Consequently
$L>AN^{1/R+s}$ and $L^3>N^{1+2/R}$, with the exact positive margins displayed.

\subsection{Diffusion inside the evolving layer and its backward factor}

For clarity, this paragraph proves a finite-layer identity independently of
all singular limits. Let $U(x,t)$ be a smooth velocity on a region and time
interval whose flow $\Phi_{t,1}(x)$ exists. Write
$J(x,t)=D_x\Phi_{t,1}(x)$ and take a smooth terminal vector field $w_1$ and a
scalar function $g(t)$. Set
\[
 G(t)=\exp\left(\int_t^1g(\tau)\,d\tau\right),\qquad
 w(x,t)=G(t)J(x,t)^{-1}w_1(\Phi_{t,1}(x)).        \tag{FS5}\label{eq:fs-flow}
\]
The inverse exists because a smooth flow has a smooth inverse. Along a
trajectory, $\Phi_{t,1}$ is constant, and differentiation of its spatial
derivative gives
$(\partial_t+U\cdot\nabla)J=-JDU$.
Differentiate $J^{-1}J=I$ to obtain
$(\partial_t+U\cdot\nabla)J^{-1}=DUJ^{-1}$.
Finally $G'=-gG$, so the product rule proves
\[
 \partial_tw+U\cdot\nabla w=w\cdot\nabla U-gw.
                                                        \tag{FS6}\label{eq:fs-damping}
\]
Thus $G\geq1$ for $g\geq0$ is the correct \emph{backward} amplification
factor for a forward damping term. Removing it would alter the terminal-value
problem.

The actual source uses the cyclic axes $e_{n+1},e_{n+2},e_{n+3}$, with indices
read modulo three, and the transverse deformation factors $g_{n2},g_{n3}$.
Its scalar damping is
\[
 g_n(t)=\nu M_n^\alpha
 \left((\partial_xg_{n2})(0,t)^2+g_{n3}(0,t)^2\right)^{\alpha/2}.
                                                        \tag{FS7}\label{eq:fs-symbol}
\]
The two phases are those in the actual transported layer, not a replacement
Cartesian Fourier mode. The scalar approximates the principal diffusion on
that layer. The force receives the residual
$\nu\Lambda^\alpha w_n-g_nw_n$.

For each use of $\Lambda^\alpha$ below, take a globally defined smooth
compactly supported vorticity on $\R^3$, as in the source's layers; the
fractional operator then has its usual Fourier-multiplier meaning. The
transport formula above applies on every region where that field has the
stated representation. Define the bilinear vector-field expression
$\mathcal B(v,w)=v\cdot\nabla w-w\cdot\nabla v$.
Equation \eqref{eq:fs-damping} gives, for every smooth $u$,
\[
 \partial_tw+\mathcal B(u,w)+\nu\Lambda^\alpha w
 =\mathcal B(u-U,w)+\nu\Lambda^\alpha w-gw.       \tag{FS8}\label{eq:fs-residual}
\]
This follows by adding and subtracting $\mathcal B(U,w)$ and using its
equation; it fixes both signs. The dissipative term has been incorporated in
the evolution, and only its approximation error appears in this formula.

The published Section 4.3.1 chooses the \emph{physical viscosity}
\[
 d=s-r>0,\quad 0<r<s,\qquad \nu=\frac{d}{200}.
                                                        \tag{FS9}\label{eq:fs-visc}
\]
It retains the factor $G_n$ omitted from the corresponding initial-layer
bound in the 2024 arXiv text. Here is the exact arithmetic of the correction.
The source's deformation and duration bounds are
\[
 |\partial_xg_{n2}(0,t)|\leq1,\quad |g_{n3}(0,t)|<2,
 \qquad
 1-T_n<34R^n A^{-R^{n-1}}\log(AN^s).
                                                        \tag{FS10}\label{eq:fs-duration}
\]
Since $5^{1/18}<2$ (raise both positive sides to power 18), these imply
\[
 0\leq g_n(t)<\frac{d}{100}M_n^\alpha.
\]
Using $t\geq T_n$, $A^{R^{n-1}}=M_n^\alpha$, and
$R^n\log(AN^s)=(a+s)\log M_n$ gives
\begin{align*}
 \log G_n(t)&=\int_t^1g_n(\tau)\,d\tau\\
 &<\frac{d}{100}M_n^\alpha(1-T_n)
 <\frac{17}{50}d(a+s)\log M_n
 <\frac25d\log M_n.
\end{align*}
Thus the source's coarser bound $G_n<M_n^{2d/5}$ follows while the factor 34
and the time duration $1-T_n$ are retained. Its additional estimate
$G_n\leq K_n$ is used in quadratic products: $G_n^2/K_n\leq G_n$.
The switching-time estimate in that section has right-hand side
\[
 G_n(T_n)K_n(T_n)^{-1}A_nM_n^r,
 \qquad K_n(T_n)=A_nM_n^s.
\]
Its exponent is therefore at most $-3d/5$, and in particular at most $-d/2$.

\subsection{Every remaining exponent and its positive summation margin}

The five error classes in published Sections 4.3.1--4.3.5 use the following
right-hand sides. We retain all logarithms when multiplying the pointwise
bound by the duration; $M=M_n$ in this subsection.

The switching term is bounded by a source constant times $M^{-3d/5}$ as just
computed. The quadratic self-interaction, including the one remaining factor
of $G_n$, has integrated scale
\[
 A_n^2L_nM^{r+2d/5-1+2\delta-\alpha}(\log M)^2
 =M^{-3d/5+2\delta}(\log M)^2.                  \tag{FS11}
\]
The equality uses $2a+\ell-1-\alpha=-s$.
The inner-velocity/outer-vorticity scale is
\[
 M^{r+2d/5+2\delta}(A/L)^{R^n}N^{R^{n-1}}\log M
 =M^{-3d/5-(a-\alpha)/2+2\delta}\log M.          \tag{FS12}
\]
Here \eqref{eq:fs-separation} gives
$a-\ell+1/R=-s-(a-\alpha)/2$.
The outer-velocity/inner-vorticity scale \eqref{eq:fs-outer} is
\[
 M^{a+r+1+2/R-(3-\delta)\ell}\log M
 =M^{-d+\ell\delta}\log M.                     \tag{FS13}
\]
This is the optimized exponent already proved above.

The diffusion estimate in published Section 4.3.5 has two terms before time
integration. Retaining the factor $G_n\leq M^{2d/5}$, their exponents are
\begin{align*}
 E_1&=\frac25d+a+\ell+\alpha+(r-1)(1-\delta)\\
    &=2\alpha-a-\frac35d+(1-r)\delta,\\
 E_2&=\frac25d+2a+\alpha+r+s+\frac2R-2\ell\\
    &=2\alpha-a-\frac35d.
\end{align*}
For the second equality use $2/R=7a$, $\ell=1+\alpha-s-2a$,
and $4s=2+3\alpha-14a$; the terms combine without discarding the $r$ dependence.
Multiplication by the duration $M^{-\alpha}\log M$ leaves the exponents
\[
 \alpha-a-\frac35d+(1-r)\delta,
 \qquad \alpha-a-\frac35d,                     \tag{FS14}
\]
with $(\log M)^2$ allowed for both terms. The source displays one logarithm
in its final integrated bound; the two logarithms directly obtained by
multiplication are kept here and have the summability proved next.

One explicit simultaneous selection is
\[
 0<\delta\leq\delta_*:=\min\left\{
 \frac d{12},\frac{a-\alpha}{4},
 \frac{1-1/R}{2},\frac1{10}\right\}.            \tag{FS15}\label{eq:fs-delta}
\]
All four entries are positive by \eqref{eq:fs-margins} and $R>1$.
The self exponent is at most $-13d/30$, and the inner exponent is no larger.
The outer exponent is smaller than $-d+(5/6)(d/12)=-67d/72$.
The first diffusion exponent is at most
$-3(a-\alpha)/4-3d/5$, since $0<r<1$ gives $(1-r)\delta\leq\delta$;
the second is negative with larger margin. In particular every displayed
exponent is strictly negative. This also ensures $M^{1/R}\leq M^{1-\delta}$,
an inequality used in the source's inner-interaction estimate.

For completeness, no logarithmic summation step is being suppressed. Given
$\epsilon>0$ and integer $m\geq0$, the function
$y^m e^{-\epsilon y/2}$ on $y\geq0$ is bounded: for $m>0$ its maximum is
$(2m/(e\epsilon))^m$, found by differentiating; for $m=0$ its bound is one.
Therefore
\[
 M_n^{-\epsilon}(\log M_n)^m
 \leq C_{m,\epsilon}N^{-\epsilon R^n/2}.
\]
Bernoulli's inequality $R^n\geq1+n(R-1)$ follows by induction from
$(1+h)^{n+1}\geq(1+nh)(1+h)\geq1+(n+1)h$ for $h=R-1>0$.
Consequently the last sequence is bounded by the summable geometric sequence
$N^{-\epsilon/2}[N^{-\epsilon(R-1)/2}]^n$. This proves summability for every
explicit scale above, including both logarithms in the diffusion term.

As an exact rational example, take
\begin{gather*}
 \alpha=\frac1{14},\quad R=2,\quad a=\frac17,
 \quad s=\frac3{56},\quad \ell=\frac{41}{56},\\
 r=\frac3{112},\quad d=\frac3{112},\quad
 \nu=\frac3{22400},\quad \delta=\frac1{448}.
\end{gather*}
Substitution in \eqref{eq:fs-parameters} proves all equalities; this $\delta$
is $d/12$ and satisfies the other three bounds in \eqref{eq:fs-delta}.
The script \path{scripts/check_forced_scaling.py} computes each exponent as an
exact rational number and verifies the identities symbolically.

\subsection{Keeping the physical viscosity and testing the classical endpoint}

For $\alpha>0$, a solution on $\R^3$ with coefficient $\nu_*>0$ has an exact spatial
transfer to any $\nu>0$. Retain
\[
 b=(\nu_*/\nu)^{1/\alpha},\quad
 \widetilde u(x,t)=b^{-1}u(bx,t),\quad
 \widetilde p(x,t)=b^{-2}p(bx,t),\quad
 \widetilde f(x,t)=b^{-1}f(bx,t).                \tag{FS16}
\]
The chain rule gives $\widetilde u_t=b^{-1}u_t$,
$\widetilde u\cdot\nabla\widetilde u=b^{-1}(u\cdot\nabla u)$,
and $\nabla\widetilde p=b^{-1}\nabla p$, all evaluated at $(bx,t)$.
The Fourier transform of $u(b\,\cdot)$ is $b^{-3}\widehat u(\xi/b)$.
Multiplication by $|\xi|^\alpha$ and inverse transformation prove
$\Lambda^\alpha\widetilde u=b^{\alpha-1}(\Lambda^\alpha u)(bx,t)$.
Since $\nu b^\alpha=\nu_*$, every term in the new equation equals
$b^{-1}$ times the old equation. Divergence vanishes by the same chain rule.
The exact $L^2$ factor for $u$ and $f$ is $b^{-5/2}$, obtained by the change of
variables $y=bx$. The supremum norm of $\nabla u$ is unchanged. A spatial
$r$-H\"older seminorm of $\nabla f$ is multiplied by $b^r$, because its
numerator is unchanged and its denominator rescales by $b^{-r}$.
The time interval is unchanged. This proves the claimed coefficient transfer
and explicitly records its data, forcing, and norm transformations.

At the classical value $\alpha=2$, the very same expressions instead give
\[
 R=\frac1{\sqrt7}<1,\qquad
 a=\frac2{\sqrt7},\qquad s=2-\sqrt7<0,
 \qquad \alpha-a=2-\frac2{\sqrt7}>0.            \tag{FS17}
\]
All three statements follow from $\sqrt7>2>1$. Thus this particular layer
choice loses increasing frequencies, positive force regularity margin, and
the negative diffusion-error exponent at once. These are exact failures of
this parametrization at $\alpha=2$. They establish no nonexistence theorem
for a different classical construction. The equations above identify the
terms whose new cancellations or estimates a reconstruction must actually
compute.

\subsection{Literal source discrepancies retained in the audit}

The saved published JATS XML has stable equation identifiers. Its
\texttt{Equ148} uses the damping sign in \eqref{eq:fs-damping}. In
\texttt{Equ192}--\texttt{Equ193}, however, the displayed force decomposition
uses $-\nu\Lambda^\alpha$ and labels a bracket ending in $-g_nw_n$ as zero.
Substitution in \eqref{eq:fs-damping} makes that bracket $-2g_nw_n$, whereas
the corrected bracket ends in $+g_nw_n$ and is zero. Identity
\eqref{eq:fs-residual} proves the corrected residual, without a sign convention
change. The definition of $K_n$ in \texttt{Equ154} omits an exponential;
the exponential is present in \texttt{Equ17} and \texttt{Equ199}. As a
stretching factor with terminal value one it must be
\[
 K_n(t)=\exp\left(\int_t^1
       -\partial_{n+2}(U_n)_{n+2}(0,0,0,\tau)\,d\tau\right).
\]
At $t=1$ the integral alone is zero, immediately contradicting that terminal
value. In \texttt{Equ197}, the duration in the first displayed exponent is
$T_n$; the integral over $[t,1]$ gives $1-t\leq1-T_n$, as calculated above.
Some powers in that XML are literally $A^n$ or $(AN^s)^n$, while the layer
definition and the arXiv TeX use $A^{R^n}$ and $(AN^s)^{R^n}$. The calculation
here explicitly fixes its scale dictionary in \eqref{eq:fs-id}; the XML and
TeX snapshots remain unchanged for inspection. These discrepancies are not
silently treated as independently validated formulas.

There is also an exact correction when time cutoffs are inserted. For a finite
family of smooth compactly supported vorticities $w_j$ on $\R^3$, set
$v_j=\mathcal K*w_j$,
$w=\sum_j\rho_jw_j$, and $u=\sum_j\rho_jv_j$, where $\rho_j$ depends only on
time. Linearity of the time derivative and of $\Lambda^\alpha$, and bilinearity
of $\mathcal B$, give
\begin{align*}
 \partial_tw+\mathcal B(u,w)+\nu\Lambda^\alpha w
 &=\sum_j\rho_j'w_j
   +\sum_j\rho_j(\partial_tw_j+\nu\Lambda^\alpha w_j)\\
 &\quad+\sum_{j,k}\rho_j\rho_k\mathcal B(v_j,w_k).
\end{align*}
In particular the self term has coefficient $\rho_j^2$, not $\rho_j$.
Their difference is $(\rho_j^2-\rho_j)\mathcal B(v_j,w_j)$. For
$0\leq\rho_j\leq1$, its coefficient has absolute value at most $1/4$, since
$\rho_j(1-\rho_j)=1/4-(\rho_j-1/2)^2$. The already displayed self-interaction
bound therefore controls this extra term with the factor $1/4$.
Ordered disjoint switching intervals make every older cutoff one while a
newer cutoff changes; this eliminates the corresponding mixed-coefficient
discrepancy. The full finite identity above remains valid without that
ordering. A locally finite family on $t<1$ satisfies the same identity on
each compact time interval because only finitely many summands are present.

The source snapshots, equation locators, exact arithmetic results, and the
independent sign derivation are retained in
\path{research/forced_program_evidence.json} and
\path{research/forced_program_sign_audit.md}. They record reconstruction
evidence and calculations, not an original discovery-session transcript.
